%Master File:lectures.tex

\lesson{Phylogenetic Trees}
\vspace{-2cm}
\begin{center}
  \includegraphics[height=10cm]{giantPanda}
\end{center}

\keywords{Trees of Life, UPGMA}
%%%%%%%%%%%%%%%%%%%%%%% Next Slide %%%%%%%%%%%%%%%%%%%%%%%
\renewcommand{\Outline}{%
\begin{slide}
\section[1]{Outline}

\begin{minipage}{8cm}
  \begin{enumerate}\squeeze
    \outlineitem{Importance of\\ Phylogeny}{bd}
    \outlineitem{Distance Based\\ Trees}{comp}
   \end{enumerate}
\end{minipage}\hfill
\begin{minipage}{14cm}
  \vfill
  \includegraphics[width=\linewidth]{giantPanda}
  \vfill
\end{minipage}
\end{slide}
\addtocounter{outlineitem}{1}
}

\setcounter{outlineitem}{1}
\Outline
\toptarget{firstoutline}
%%%%%%%%%%%%%%%%%%%%%%% Next Slide %%%%%%%%%%%%%%%%%%%%%%%

\begin{slide}
\section{Giant Panda}

\begin{PauseHighLight}
  \begin{itemize}
  \item Is a giant panda a bear or a raccoon?\pause
  \item When they were first discovered in Europe scientists couldn't
    decide---they had features of both\pause
  \item They were named after red pandas that are raccoons\pause
  \item This was only resolved in 1985 when Steven O'Brien and
    colleagues used sequence matching from 500\,000 nucleotides to
    finally produce a definitive evolutionary tree\pause
  \end{itemize}
\end{PauseHighLight}

\end{slide}

%%%%%%%%%%%%%%%%%%%%%%% Next Slide %%%%%%%%%%%%%%%%%%%%%%%

\begin{slide}
\section[-2]{The Tree of Bears}

\begin{center}
  \includegraphics[width=0.9\linewidth]{giantPanda}
\end{center}
\end{slide}

%%%%%%%%%%%%%%%%%%%%%%% Next Slide %%%%%%%%%%%%%%%%%%%%%%%

\begin{slide}
\section{Out of Africa}

\begin{PauseHighLight}
  \begin{itemize}
  \item In the mid 1980s Allan Wilson and colleagues studied the
    evolutionary tree of human mitrochondria and decided that all races
    outside Africa shared a common ancestor who lived in Africa 200\,000
    years ago\pause
  \item Hominids had left Africa much earlier (e.g. Neanderthals) so
    this theory postulated that modern humans would have displaced all
    existing hominids outside of Africa\pause
  \end{itemize}
\end{PauseHighLight}

\end{slide}

%%%%%%%%%%%%%%%%%%%%%%% Next Slide %%%%%%%%%%%%%%%%%%%%%%%

\begin{slide}
\section{Modern Humans}

\begin{center}
  \includegraphics[width=\linewidth]{OutOfAfrica}\\
  \tiny{By NordNordWest - File:Spreading homo sapiens ru.svg by Urutseg, Public Domain, https://commons.wikimedia.org/w/index.php?curid=34697001}
\end{center}
\end{slide}

%%%%%%%%%%%%%%%%%%%%%%% Next Slide %%%%%%%%%%%%%%%%%%%%%%%

\begin{slide}
\section{Mitrochondrial DNA}

\begin{PauseHighLight}
  \begin{itemize}
  \item Mitrochondria are the powerhouse of cells\pause
  \item They would have originally been an independent organism that
    somehow got captured by eukaryotic cells\pause
  \item They contain their own DNA\pause
  \item They are inherited through the mother's ova\pause
  \item They therefore reproduce asexually\pause
  \end{itemize}
\end{PauseHighLight}

\end{slide}

%%%%%%%%%%%%%%%%%%%%%%% Next Slide %%%%%%%%%%%%%%%%%%%%%%%

\begin{slide}
\section{Asexual Inheritance}

\begin{PauseHighLight}
  \begin{itemize}
  \item Asexual inheritance is much easier to study\pause
  \item Populations will share a common ancestor because of
    \begin{itemize}
    \item Natural coalescence\pause
    \item Going through a bottle neck (e.g. leaving Africa)\pause
    \end{itemize}
  \item Natural coalescence happens by chance, but will be enhanced by
    selection\pause
  \end{itemize}
\end{PauseHighLight}

\end{slide}

%%%%%%%%%%%%%%%%%%%%%%% Next Slide %%%%%%%%%%%%%%%%%%%%%%%

\begin{slide}
\section[-2]{Coalescence}

\pause
\pb
\pauselevel{=1}
\begin{center}
  \multipdf[width=\linewidth]{coalescence}\pause
\end{center}
\begin{itemize}
\item If you go far enough back in time an asexual population will have a
  common ancestor\pause
\end{itemize}
\end{slide}

%%%%%%%%%%%%%%%%%%%%%%% Next Slide %%%%%%%%%%%%%%%%%%%%%%%

\begin{slide}
\section{Cross-Breeding}

\begin{PauseHighLight}
  \begin{itemize}
  \item Didn't modern man cross-breed with other hominids?\pause
  \item This was unknown even 10 years ago\pause, but of course they
    did\pauseb
  \item However, the numbers were not large and due to coalescence (the
    coalescence time being small because there was not huge
    cross-breeding) all traces of Neanderthal  mitrochondrial DNA vanished\pause
  \item The story is very different for the rest of our DNA\pause
  \item Normal DNA evolves sexually\pause
  \end{itemize}
\end{PauseHighLight}

\end{slide}

%%%%%%%%%%%%%%%%%%%%%%% Next Slide %%%%%%%%%%%%%%%%%%%%%%%

\begin{slide}
\section{The Selfish Gene}

\begin{PauseHighLight}
  \begin{itemize}
  \item With crossover you cannot build a tree of life as you inherit
    from two parents\pause
  \item However, each gene (at least if it doesn't get disrupted by
    crossover) can be mapped to its own ancestral tree\pause
  \item Because there are lots of genes it is much more likely for genes
    from different origins to survive (particularly if they have
    different fitness)\pause
  \item In fact, we carry DNA that clearly came from Neanderthal and a
    few other early hominids\pause
  \item These genes seem to have been selected for\pause
  \end{itemize}
\end{PauseHighLight}

\end{slide}

%%%%%%%%%%%%%%%%%%%%%%% Next Slide %%%%%%%%%%%%%%%%%%%%%%%

\begin{slide}
\section{Horizontal Gene Transfer}

\begin{minipage}{0.52\linewidth}
\begin{PauseHighLight}
  \begin{itemize}
  \item Although bacteria and archaea reproduce asexually they regularly
    exchange plasmids\pause
  \item These are loops of DNA\pause
  \item This is the primary cause of\\ antibiotic resistance\pause
  \item The tree of life is not quite a tree\pause
  \end{itemize}
\end{PauseHighLight}  
\end{minipage}\hfil
\begin{minipage}{0.45\linewidth}
  \begin{center}
    \includegraphics[width=0.9\linewidth]{Horizontal-gene-transfer}\\
    \tiny{By Dr. Smets and perhaps others - Barth F. Smets, Ph.D., with
      permission, Attribution,
      https://commons.wikimedia.org/w/index.php?curid=5720886}    
  \end{center}
\end{minipage}

\end{slide}

%%%%%%%%%%%%%%%%%%%%%%% Next Slide %%%%%%%%%%%%%%%%%%%%%%%

\begin{slide}
\section{Ligers and Tigons}

\begin{PauseHighLight}
  \begin{itemize}
  \item Most higher animals form species where individuals between
    species don't inter-breed\pause{} very much\pauseb
  \item In captivity ligers are crosses between male lions and female
    tigers  and tigons are crosses between male tigers and female
    lions\pause
  \item In the wild lions and tigers no longer overlap, but in the past
    there was some report of ligers (which are typically enormous)\pause
  \item Like most cross-breeds they have reduced fertility (one sex is
    often inviable or sterile) so that cross breeds die out\pause
  \item Because of this lack of significant cross-breeding we can model
    most of evolution at a species level as a tree\pause
  \end{itemize}
\end{PauseHighLight}

\end{slide}

%%%%%%%%%%%%%%%%%%%%%%% Next Slide %%%%%%%%%%%%%%%%%%%%%%%

\begin{slide}
\section{Evolution and HIV}

\begin{PauseHighLight}
  \begin{itemize}
  \item We often think of evolution as happening over geological
    time\pause
  \item But the HIV virus evolved over decades from a disease that
    infected African primates to one that infects man\pause
  \item Furthermore it evolves over months and years to evade our
    immune system\pause
  \item Tracking these changes can be important in understanding and
    controlling the disease\pause
  \item The virus evolves asexually so we can build an evolution
    tree\pause
  \end{itemize}
\end{PauseHighLight}

\end{slide}



%%%%%%%%%%%%%%%%%%%%%%% Next Slide %%%%%%%%%%%%%%%%%%%%%%%

\begin{slide}
\section{Evolution of Cancer}

\begin{PauseHighLight}
  \begin{itemize}
  \item Cells even in higher organisms predominantly reproduce asexually
    (they continuously divide)\pause
  \item Over time the genome can get disrupted which if we are unlucky
    can grow into tumours\pause
  \item Cancer cells accumulate mutations, some of which will die out
    others of which will replicate faster that other cells\pause
  \item Understanding the development of cancer is often crucial to its
    treatment\pause
  \end{itemize}
\end{PauseHighLight}

\end{slide}

%%%%%%%%%%%%%%%%%%%%%%% Next Slide %%%%%%%%%%%%%%%%%%%%%%%

\begin{slide}
\section{Orthologues and Paralogues}

\begin{PauseHighLight}
  \begin{itemize}
  \item Genes that diverge due to speciation are known as
    orthologues\pause
  \item Genes that diverge due to gene duplication are known as
    paralogues\pause
  \item If we are interested in understanding the relationship between
    species (e.g. pandas) we have to make sure we choose
    orthologues\pause
  \item But if we are interested in the different types of haemoglobins
    then we have to look at paralogues\pause
  \end{itemize}
\end{PauseHighLight}


\end{slide}

%%%%%%%%%%%%%%%%%%%%%%% Next Slide %%%%%%%%%%%%%%%%%%%%%%%

\begin{slide}
\section[-1]{}

\begin{center}
  \includegraphics[width=0.9\linewidth]{haemoglobins1}
\end{center}
\end{slide}

%%%%%%%%%%%%%%%%%%%%%%% Next Slide %%%%%%%%%%%%%%%%%%%%%%%

\begin{slide}
\section[-2]{}


\begin{center}
  \includegraphics[width=0.6\linewidth]{haemoglobins}
\end{center}
\end{slide}


%%%%%%%%%%%%%%%%%%%%%%% Next Slide %%%%%%%%%%%%%%%%%%%%%%%

\begin{slide}
\section{Three Camps}

\begin{PauseHighLight}
  \begin{itemize}
  \item Although you would think that evolutionary trees should be
    pretty well understood, there is a huge debate about how to do this
    right\pause
  \item There are three bitter rival camps in the field
    \begin{itemize}
    \item Evolutionary Taxonomy
    \item Phenetics or numerical taxonomy
    \item Cladistics\pause
    \end{itemize}
  \item If you are interested you can delve in the philosophical
    differences, but we will ignore them\pause
  \end{itemize}
\end{PauseHighLight}

\end{slide}


%%%%%%%%%%%%%%%%%%%%%%% Next Slide %%%%%%%%%%%%%%%%%%%%%%%
\Outline % Ultrametric
%%%%%%%%%%%%%%%%%%%%%%% Next Slide %%%%%%%%%%%%%%%%%%%%%%%

\begin{slide}
\section{Building Phylogenetic Trees}

\begin{PauseHighLight}
  \begin{itemize}
  \item Building evolutionary trees (also known as phylogenetic trees) based
    on some measure of evolutionary distance is scientifically
    important\pause
  \item Traditionally morphological features where used, but today
    sequence edit distance is usually used\pause
  \item There are often many plausible tree that can explain the
    data\pause
  \item Efficient methods that find the best evolutionary trees still
    don't exist---current algorithms are compromises\pause
  \end{itemize}
\end{PauseHighLight}

\end{slide}

%%%%%%%%%%%%%%%%%%%%%%% Next Slide %%%%%%%%%%%%%%%%%%%%%%%

\begin{slide}
\section{The Molecular Clock}

\begin{PauseHighLight}
  \begin{itemize}
  \item It is a nice picture to see the accumulation of mutations acts as a
    clock telling us when species separated\pause
  \item Unfortunately, mutations are random and strongly affected by
    selection pressure\pause
  \item Also mutations are often very complex, involve a duplication of
    a stretch of DNA\pause
  \item A lot of care is needed to use mutations as a molecular clock,
    including using as large sequences as possible\pause
  \end{itemize}
\end{PauseHighLight}

\end{slide}

%%%%%%%%%%%%%%%%%%%%%%% Next Slide %%%%%%%%%%%%%%%%%%%%%%%

\begin{slide}
\section{Two Types of Trees}

\begin{PauseHighLight}
  \begin{itemize}
  \item There are two main types of evolutionary trees
    \begin{itemize}
    \item \emph{Distance-based trees}, where we have a table of
      ``distances'' between species and we wish to find a tree
      explaining these distances\pause
    \item \emph{Sequence based trees}, where we are given a set of
      sequences and we want to find a tree with a minimum number of
      mutations per link (maximum parsimony) that explains the
      data\pause
    \end{itemize}
  \end{itemize}
\end{PauseHighLight}

\end{slide}

%%%%%%%%%%%%%%%%%%%%%%% Next Slide %%%%%%%%%%%%%%%%%%%%%%%

\begin{slide}
\section[-2]{Ultrametric Trees}

\begin{PauseHighLight}
  \begin{itemize}
  \item We consider binary trees\pause
  \item If we had a perfect molecular clock then evolutionary trees
    to existing species will have an ultrametric structure\pause
    \begin{center}
      \includegraphics[width=0.9\linewidth]{ultrametric}
    \end{center}
  \item For any three nodes $i$, $j$ and $k$ with distances $d_{ij}$,
    $d_{ik}$ and $d_{jk}$ two distances will be identical, and one will
    be smaller than these\pause
  \end{itemize}
\end{PauseHighLight}

\end{slide}

%%%%%%%%%%%%%%%%%%%%%%% Next Slide %%%%%%%%%%%%%%%%%%%%%%%

\begin{slide}
\section{Real Trees}

\begin{PauseHighLight}
  \begin{itemize}
  \item Because we cannot accurately measure divergence time, in practice
    it is very unlikely that the table of distances can be modelled with
    an ultrametric tree\pause
  \item We can however build an ultrametric tree through a clustering
    algorithm known as \emph{Unweighted Pair Group Method using
      Arithmetic averaging} UPGMA\pause
  \item This is an easy algorithm where we build a tree from the leaves
    upwards\pause
  \item We select the two closest subtrees remove all the nodes and
    replace the two joined subtree with a new subtree\pause
  \end{itemize}
\end{PauseHighLight}

\end{slide}

%%%%%%%%%%%%%%%%%%%%%%% Next Slide %%%%%%%%%%%%%%%%%%%%%%%

\begin{slide}
\section{UPGMA}
\pause
\pb
\pauselevel{=1}
\begin{center}
  \multipdf[width=\linewidth]{upgma}\pause
\end{center}
\end{slide}

%%%%%%%%%%%%%%%%%%%%%%% Next Slide %%%%%%%%%%%%%%%%%%%%%%%

\begin{slide}
\section[-2]{Distance Calculation for UPGMA}

\begin{PauseHighLight}
  \begin{itemize}
  \item The UPGMA algorithm starts from a table of distances\pause
  \item The distance between two clusters $\mathcal{C}_k$ and
    $\mathcal{C}_l$ is given by
    \begin{align*}
      d_{kl} = \frac{1}{|\mathcal{C}_k|\,|\mathcal{C}_l|}
      \sum_{\substack{p\in\mathcal{C}_k\\q\in\mathcal{C}_l}} d_{pq}\pause
    \end{align*}
  \item If we join two clusters so that $\mathcal{C}_k = \mathcal{C}_i
    \cup \mathcal{C}_j$ then the distance to a cluster $\mathcal{C}_l$
    is
    \begin{align*}
      d_{kl} = \frac{d_{il}\,|\mathcal{C}_i| + d_{jl}\,|\mathcal{C}_j|}
      {|\mathcal{C}_i| + |\mathcal{C}_j|}\pause
    \end{align*}
  \item If the distances are not ultrametric then UPGMA can give poor
    trees\pause
  \end{itemize}
\end{PauseHighLight}

\end{slide}

%%%%%%%%%%%%%%%%%%%%%%% Next Slide %%%%%%%%%%%%%%%%%%%%%%%

\begin{slide}
\section[-2]{Weakness of UPGMA}
\pause
\pb
\pauselevel{=1}
\begin{itemize}
\item Note that organism 5 and 7 are quite close, but UPGMA distorts
  this\pauseh 
\end{itemize}
\pauselevel{=1}
\begin{center}
  \multipdf[width=0.9\linewidth]{upgma}\pause
\end{center}

\end{slide}

%%%%%%%%%%%%%%%%%%%%%%% Next Slide %%%%%%%%%%%%%%%%%%%%%%%

\begin{slide}
\section{Reflection on Ultrametricity}

\begin{PauseHighLight}
  \begin{itemize}
  \item UPGMA imposes an ultrametric structure on the derived
    evolutionary trees\pause
  \item If mutations accumulated at a constant rate then evolutionary
    trees should have an ultrametric structure\pause
  \item Unfortunately mutations experience selection and in some cases
    they are very non-trivial (genome copy mutations)\pause
  \item There are stretches of ``junk'' DNA where mutations appear to
    accumulate at a reasonable constant rate, but these are not well
    studied or stable regions\pause
  \end{itemize}
\end{PauseHighLight}

\end{slide}


%%%%%%%%%%%%%%%%%%%%%%% Next Slide %%%%%%%%%%%%%%%%%%%%%%%

\begin{slide}
\section{Summary}

\begin{PauseHighLight}
  \begin{itemize}
  \item Evolutionary trees provide very important scientific insights\pause
  \item Building evolutionary trees is non-trivial\pause
  \item If the distances are approximately ultrametric then UPGMA will
    generate a good tree\pause
  \item Need to look at other ways of building evolutionary trees\pause
  \end{itemize}
\end{PauseHighLight}

\end{slide}

